\uuid{KtZu}
\titre{Degré}
\auteur{bodin ; hanani ; mzari}
\organisation{exo7}
\licence{CC BY-NC-SA 4.0}
\datecreate{2020-01-20}
\niveau{L1}
\module{Algèbre}
\chapitre{Polynôme, fraction rationnelle}
\sousChapitre{Définition, degré, produit}
\difficulte{1}
\theme{Polynômes, degré}
\exo7qcmid{lille-1.polynomes.003}

\contenu{
\begin{question}
\qclassification{105.05}
\qtags{motcle=degré}
Soient $P(X)$ et $Q(X)$ deux polynômes unitaires de degré $n\ge1$.
Quelles sont les assertions vraies ?
\begin{answers}
    \good{$P+Q$ est un polynôme de degré $n$.}

    \bad{$P-Q$ est un polynôme de degré $n$.}

    \good{$P \times Q$ est un polynôme de degré $n+n=2n$.}

    \bad{$P/Q$ est un polynôme de degré $n-n=0$.}
\end{answers}
\begin{explanations}
Le quotient de deux polynômes n'est pas un polynôme.
\end{explanations}
\end{question}
}
